% =============================================================================
% Section 2: Background and Related Work
% =============================================================================
\section{Background and Related Work}
\label{sec:background}

% ---------------------------------------------------------------------------
% 2.1 Foundation Models for Brain MRI
% ---------------------------------------------------------------------------
\subsection{Foundation Models for Brain MRI}
\label{sec:bg:foundation}

The paradigm of pretraining large encoder--decoder architectures on massive,
diverse datasets and subsequently fine-tuning on downstream tasks has
transformed medical image analysis \cite{he2023foundation,moor2023foundation}.
BrainSegFounder \cite{cox2024brainsegfounder} instantiates this paradigm for
3D brain MRI segmentation through a three-stage training pipeline:

\begin{enumerate}[leftmargin=*]
  \item \textbf{Self-supervised pretraining.} A SwinUNETR encoder
    \cite{hatamizadeh2022swinunetr,tang2022self} is pretrained on 41,400+
    brain MRI subjects from the UK Biobank using a combination of masked image
    modelling, contrastive learning, and rotation prediction. This stage uses a
    $96^3$ voxel region-of-interest (ROI) and produces general-purpose
    anatomical features.

  \item \textbf{Supervised pretraining.} The encoder is paired with the
    SwinUNETR decoder and trained on the BraTS~2021 glioma segmentation
    benchmark \cite{baid2021rsna} at $128^3$ ROI, producing tumour-specific
    features.

  \item \textbf{Per-fold fine-tuning.} Final fine-tuning on BraTS~2021 with
    cross-validation yields the released checkpoints.
\end{enumerate}

\subsubsection{SwinUNETR Architecture}

SwinUNETR combines the Swin Transformer \cite{liu2021swin} with a
UNet-style encoder--decoder structure. The encoder consists of a patch
embedding layer followed by four transformer stages with shifted window
self-attention. For BSF-Tiny (the variant used in this work), the
architectural parameters are:

\begin{table}[H]
  \centering
  \caption{SwinUNETR encoder (BSF-Tiny) architecture. Input: $[B, 4, 128^3]$.}
  \label{tab:swinunetr_arch}
  \begin{tabular}{lcccc}
    \toprule
    \textbf{Stage} & \textbf{Output shape} & \textbf{Channels} & \textbf{Heads} & \textbf{Blocks} \\
    \midrule
    Patch embed & $[B, 48, 64^3]$  & 48  & --- & --- \\
    Stage 1     & $[B, 96, 32^3]$  & 96  & 3   & 2 \\
    Stage 2     & $[B, 192, 16^3]$ & 192 & 6   & 2 \\
    Stage 3     & $[B, 384, 8^3]$  & 384 & 12  & 2 \\
    Stage 4     & $[B, 768, 4^3]$  & 768 & 24  & 2 \\
    encoder10   & $[B, 768, 4^3]$  & 768 & --- & --- \\
    \bottomrule
  \end{tabular}
\end{table}

The bottleneck layer \texttt{encoder10} applies a residual convolutional
block preserving the $768$-channel, $4^3$-spatial resolution from Stage~4.
Global Average Pooling (GAP) over the spatial dimensions produces the feature
vector $\mathbf{h} \in \R^{768}$ used throughout our pipeline.

The segmentation decoder mirrors the encoder with skip connections, producing
a 3-channel sigmoid output following the BraTS convention: Channel~0 encodes
Tumour Core (TC), Channel~1 encodes Whole Tumour (WT), and Channel~2 encodes
Enhancing Tumour (ET), using hierarchical overlapping labels rather than
mutually exclusive classes.

% ---------------------------------------------------------------------------
% 2.2 Low-Rank Adaptation (LoRA)
% ---------------------------------------------------------------------------
\subsection{Low-Rank Adaptation}
\label{sec:bg:lora}

Low-Rank Adaptation (LoRA) \cite{hu2022lora} provides parameter-efficient
fine-tuning by injecting trainable rank-$r$ decompositions into frozen
pretrained weight matrices. For a pretrained weight matrix
$\mathbf{W}_0 \in \R^{d \times k}$, LoRA parameterises the adapted weight as:
\begin{equation}
  \mathbf{W} = \mathbf{W}_0 + \frac{\alpha}{r}\, \mathbf{B}\mathbf{A},
  \quad \mathbf{B} \in \R^{d \times r},\; \mathbf{A} \in \R^{r \times k},
  \label{eq:lora}
\end{equation}
where $r \ll \min(d, k)$ is the rank, $\alpha$ is a scaling hyperparameter,
$\mathbf{A}$ is initialised from $\mathcal{N}(0, \sigma^2)$, and
$\mathbf{B} = \mathbf{0}$ so that the adapted model initially reproduces the
pretrained model exactly. The number of trainable parameters per adapted layer
is $r(d + k)$, compared to $dk$ for full fine-tuning.

The theoretical foundation rests on the observation by Aghajanyan et al.\
\cite{aghajanyan2021intrinsic} that pretrained language models have a low
intrinsic dimensionality: the parameter space can be projected to a
low-dimensional subspace without significant loss in fine-tuning performance.
LoRA operationalises this insight by constraining the weight update to lie in a
rank-$r$ subspace.

\textbf{DoRA} \cite{liu2024dora} extends LoRA by decomposing the weight
update into magnitude and direction components:
\begin{equation}
  \mathbf{W} = m \cdot \frac{\mathbf{W}_0 + \mathbf{B}\mathbf{A}}
    {\|\mathbf{W}_0 + \mathbf{B}\mathbf{A}\|_c},
  \label{eq:dora}
\end{equation}
where $m$ is a learnable magnitude vector and $\|\cdot\|_c$ denotes
column-wise normalisation. This decomposition was shown to more closely
approximate full fine-tuning behaviour in language models; we evaluate it as
an ablation condition (A2, \cref{sec:eval:ablations}).

% ---------------------------------------------------------------------------
% 2.3 Supervised Disentangled Projection
% ---------------------------------------------------------------------------
\subsection{Disentangled Representation Learning}
\label{sec:bg:disentanglement}

Disentangled representations encode independent generative factors of
variation in separate latent dimensions or subspaces. Higgins et al.\
\cite{higgins2018towards} formalised this via group theory: a representation
is disentangled with respect to a group decomposition
$G = G_1 \times \cdots \times G_n$ if there exists a corresponding
decomposition $Z = Z_1 \oplus \cdots \oplus Z_n$ where each $Z_i$ is affected
only by $G_i$.

The impossibility theorem of Locatello et al.\ \cite{locatello2019challenging}
established that unsupervised disentanglement is theoretically impossible
without inductive biases: for any generative model producing disentangled
representations, there exists another model with identical marginal
distribution but entangled latents. Their follow-up work
\cite{locatello2020disentangling} demonstrated that as few as 100 labelled
examples suffice to resolve this impossibility through supervision.

These results motivate our Supervised Disentangled Projection (SDP) approach:
rather than relying on VAE-based unsupervised disentanglement (which we
empirically found to suffer from posterior collapse and KL distortion), we
use explicit semantic supervision from segmentation-derived features to
anchor each latent partition to a known generative factor.

\subsubsection{VICReg Regularisation}

VICReg \cite{bardes2022vicreg} defines three regularisation terms for
representation learning. The \textit{variance} term prevents dimensional
collapse via a hinge loss on per-dimension standard deviation:
\begin{equation}
  \mathcal{L}_{\text{var}}(\mathbf{Z}) = \frac{1}{d}\sum_{j=1}^{d}
    \max\!\bigl(0,\; \gamma - \sqrt{\operatorname{Var}(z^{(j)}) + \epsilon}
    \,\bigr),
  \label{eq:vicreg_var}
\end{equation}
where $\gamma = 1$ is the target standard deviation. The \textit{covariance}
term decorrelates dimensions by penalising off-diagonal elements of the batch
covariance matrix $\mathbf{C}$:
\begin{equation}
  \mathcal{L}_{\text{cov}}(\mathbf{Z}) = \frac{1}{d}\sum_{i \neq j}
    C_{ij}^2.
  \label{eq:vicreg_cov}
\end{equation}

\subsubsection{Distance Correlation}

VICReg's covariance penalty captures only linear dependencies. Distance
correlation \cite{szekely2007measuring} detects all forms of statistical
dependence between random vectors $X$ and $Y$ of potentially different
dimensions:
\begin{equation}
  \dCor(X, Y) = \frac{\dCov(X, Y)}
    {\sqrt{\dVar(X) \cdot \dVar(Y)}},
  \label{eq:dcor}
\end{equation}
where $\dCov$ is computed from doubly-centred pairwise distance matrices.
Critically, $\dCor(X, Y) = 0$ if and only if $X \perp\!\!\!\perp Y$ (for
random vectors with finite first moments), providing a necessary and
sufficient condition for statistical independence.

\subsubsection{Spectral Normalisation}

Spectral normalisation \cite{miyato2018spectral} constrains each linear layer
$\mathbf{W}$ by dividing by its largest singular value
$\sigma(\mathbf{W})$, yielding a weight matrix with spectral norm exactly~1.
For an $L$-layer MLP with 1-Lipschitz activations (e.g., ReLU, GELU), this
bounds the global Lipschitz constant: $\operatorname{Lip}(f) \leq
\prod_{\ell=1}^{L} 1 = 1$. The contraction property
$\|\mathbf{z}_1 - \mathbf{z}_2\| \leq \|\mathbf{h}_1 - \mathbf{h}_2\|$
ensures that encoder feature noise is not amplified through the projection,
which is essential for stable downstream GP prediction
\cite{gouk2021regularisation}.

% ---------------------------------------------------------------------------
% 2.4 Gaussian Process Regression
% ---------------------------------------------------------------------------
\subsection{Gaussian Process Regression}
\label{sec:bg:gp}

A Gaussian process (GP) is a distribution over functions such that any finite
collection of function values is jointly Gaussian
\cite{rasmussen2006gaussian}:
\begin{equation}
  f(\mathbf{t}) \sim \GP\bigl(m(\mathbf{t}),\; k(\mathbf{t}, \mathbf{t}')\bigr),
  \label{eq:gp_def}
\end{equation}
where $m(\mathbf{t})$ is the mean function and $k(\mathbf{t}, \mathbf{t}')$
is the covariance (kernel) function. Given $n$ observations
$\mathbf{y} = f(\mathbf{t}) + \boldsymbol{\epsilon}$ with
$\boldsymbol{\epsilon} \sim \mathcal{N}(\mathbf{0}, \sigma_n^2 \mathbf{I})$,
the predictive distribution at test inputs $\mathbf{t}_*$ is:
\begin{align}
  \mathbf{f}_* | \mathbf{t}_*, \mathbf{t}, \mathbf{y}
    &\sim \mathcal{N}(\boldsymbol{\mu}_*, \boldsymbol{\Sigma}_*), \\
  \boldsymbol{\mu}_* &= m(\mathbf{t}_*) + \mathbf{K}_*^\top
    (\mathbf{K} + \sigma_n^2\mathbf{I})^{-1}(\mathbf{y} - m(\mathbf{t})),
    \label{eq:gp_mean} \\
  \boldsymbol{\Sigma}_* &= \mathbf{K}_{**} - \mathbf{K}_*^\top
    (\mathbf{K} + \sigma_n^2\mathbf{I})^{-1}\mathbf{K}_*,
    \label{eq:gp_var}
\end{align}
where $\mathbf{K} = k(\mathbf{t}, \mathbf{t})$,
$\mathbf{K}_* = k(\mathbf{t}, \mathbf{t}_*)$, and
$\mathbf{K}_{**} = k(\mathbf{t}_*, \mathbf{t}_*)$.

\subsubsection{Mat\'{e}rn Kernels}

The Mat\'{e}rn family of kernels provides control over the smoothness of
the GP sample paths via the parameter $\nu$:
\begin{equation}
  k_{\nu}(\Delta t) = \sigma_f^2 \frac{2^{1-\nu}}{\Gamma(\nu)}
    \left(\frac{\sqrt{2\nu}\,|\Delta t|}{\ell}\right)^\nu
    K_\nu\!\left(\frac{\sqrt{2\nu}\,|\Delta t|}{\ell}\right),
  \label{eq:matern}
\end{equation}
where $\ell$ is the length-scale, $\sigma_f^2$ the signal variance, and
$K_\nu$ the modified Bessel function. For $\nu = 5/2$, the kernel yields
twice-differentiable sample paths:
\begin{equation}
  k_{5/2}(\Delta t) = \sigma_f^2
    \left(1 + \frac{\sqrt{5}\,|\Delta t|}{\ell}
    + \frac{5\,\Delta t^2}{3\ell^2}\right)
    \exp\!\left(-\frac{\sqrt{5}\,|\Delta t|}{\ell}\right).
  \label{eq:matern52}
\end{equation}
This is our default temporal kernel for volume dynamics, encoding the prior
belief that tumour volume evolves smoothly over time.

\subsubsection{Multi-Output GPs and the Intrinsic Coregionalisation Model}

Multi-Output Gaussian Processes (MOGPs) extend scalar GPs to vector-valued
functions $\mathbf{f}: \R \to \R^D$, capturing correlations across output
dimensions. The Intrinsic Coregionalisation Model (ICM) \cite{alvarez2012kernels}
parameterises the cross-output covariance as a Kronecker product:
\begin{equation}
  \mathbf{K}(\mathbf{t}, \mathbf{t}') =
    \mathbf{B} \otimes k_t(\mathbf{t}, \mathbf{t}'),
  \label{eq:icm}
\end{equation}
where $\mathbf{B} \in \R^{D \times D}$ is a positive semi-definite
coregionalisation matrix encoding inter-output correlations, and
$k_t$ is a scalar temporal kernel shared across outputs. The Linear Model of
Coregionalisation (LMC) generalises this to a sum over $Q$ latent functions:
\begin{equation}
  \mathbf{K}(\mathbf{t}, \mathbf{t}') =
    \sum_{q=1}^{Q} \mathbf{B}_q \otimes k_q(\mathbf{t}, \mathbf{t}').
  \label{eq:lmc}
\end{equation}

% ---------------------------------------------------------------------------
% 2.5 Related Work
% ---------------------------------------------------------------------------
\subsection{Related Work}
\label{sec:bg:related}

\subsubsection{Brain Tumour Growth Prediction}

Tumour growth prediction from longitudinal MRI has been addressed through
several paradigms. Classical approaches fit parametric growth models
(exponential, Gompertz, logistic) to volume trajectories
\cite{nakasu2005serial,hashimoto2012natural}. Reaction--diffusion models
\cite{swanson2000quantitative,hormuth2021image} incorporate spatial diffusion
but require tissue-specific parameterisation and are computationally expensive.

Recent deep learning approaches include spatio-temporal convolutional LSTMs
\cite{zhang2020spatio} achieving 83.2\% Dice on glioma growth prediction,
GP-GAN \cite{elazab2020gp} combining Gaussian processes with GANs, and neural
fields conditioned on time \cite{chen2024vestibular} for vestibular schwannoma
growth. Puglisi et al.\ \cite{puglisi2025brain} demonstrated that a linear
mixed-effects model in a learned latent space produces clinically meaningful
progression predictions for brain atrophy, trained on 11,730~MRIs.

For meningiomas specifically, most existing work is limited to retrospective
volumetric analysis \cite{hashimoto2012natural,ius2020management,
nakamura2003natural}. To our knowledge, no prior work has applied foundation
model-derived latent spaces with GP-based temporal modelling to meningioma
growth prediction.

\subsubsection{Latent Space Temporal Prediction}

The Joint Embedding Predictive Architecture (JEPA) framework
\cite{lecun2022path} advocates predicting in representation space rather than
pixel space, avoiding the waste of capacity on irrelevant details. I-JEPA
\cite{assran2023self} and V-JEPA \cite{bardes2024revisiting} validated this
for images and video respectively. Longitudinal Self-Supervised Learning
(LSSL) \cite{zhao2021longitudinal} disentangles brain-age factors from latent
MRI representations on 811 ADNI subjects. Our pipeline aligns with this
``predict in latent space'' principle while using GPs instead of neural
predictors, providing closed-form uncertainty quantification for clinical
decision support.

\subsubsection{Domain Adaptation in Medical Imaging}

Domain adaptation for medical image analysis has been addressed through
adversarial training \cite{kamnitsas2017unsupervised}, style transfer
\cite{chen2022maxstyle}, and more recently through parameter-efficient
fine-tuning of foundation models \cite{dutt2023parameter}. Ben-David et al.\
\cite{bendavid2010theory} established the theoretical framework relating
domain divergence to target risk. Our work contributes empirical evidence on
LoRA adaptation for 3D medical imaging foundation models, including
dual-domain training that attempts to maintain glioma competence while
acquiring meningioma specialisation.
